% SNN 本地学习的逐层信息保留与损失定位（2026-10）
\documentclass[journal,12pt,onecolumn]{IEEEtran}
\IEEEoverridecommandlockouts
\usepackage{cite}
\usepackage{amsmath, amssymb, bm, graphicx, enumitem}
\usepackage{algorithmic}
\usepackage{algorithm}
\usepackage{booktabs}
\usepackage{textcomp}
\usepackage{xcolor}
\usepackage{tabularx}
\usepackage{makecell}
\usepackage{tikz}
\usetikzlibrary{arrows.meta, positioning}
\usepackage[UTF8]{ctex}

%%%%% NEW MATH DEFINITIONS %%%%%

\usepackage{amsmath,amsfonts,bm}

% Mark sections of captions for referring to divisions of figures
\newcommand{\figleft}{{\em (Left)}}
\newcommand{\figcenter}{{\em (Center)}}
\newcommand{\figright}{{\em (Right)}}
\newcommand{\figtop}{{\em (Top)}}
\newcommand{\figbottom}{{\em (Bottom)}}
\newcommand{\captiona}{{\em (a)}}
\newcommand{\captionb}{{\em (b)}}
\newcommand{\captionc}{{\em (c)}}
\newcommand{\captiond}{{\em (d)}}

% Highlight a newly defined term
\newcommand{\newterm}[1]{{\bf #1}}


% Figure reference, lower-case.
\def\figref#1{figure~\ref{#1}}
% Figure reference, capital. For start of sentence
\def\Figref#1{Figure~\ref{#1}}
\def\twofigref#1#2{figures \ref{#1} and \ref{#2}}
\def\quadfigref#1#2#3#4{figures \ref{#1}, \ref{#2}, \ref{#3} and \ref{#4}}
% Section reference, lower-case.
\def\secref#1{section~\ref{#1}}
% Section reference, capital.
\def\Secref#1{Section~\ref{#1}}
% Reference to two sections.
\def\twosecrefs#1#2{sections \ref{#1} and \ref{#2}}
% Reference to three sections.
\def\secrefs#1#2#3{sections \ref{#1}, \ref{#2} and \ref{#3}}
% Reference to an equation, lower-case.
\def\eqref#1{equation~\ref{#1}}
% Reference to an equation, upper case
\def\Eqref#1{Equation~\ref{#1}}
% A raw reference to an equation---avoid using if possible
\def\plaineqref#1{\ref{#1}}
% Reference to a chapter, lower-case.
\def\chapref#1{chapter~\ref{#1}}
% Reference to an equation, upper case.
\def\Chapref#1{Chapter~\ref{#1}}
% Reference to a range of chapters
\def\rangechapref#1#2{chapters\ref{#1}--\ref{#2}}
% Reference to an algorithm, lower-case.
\def\algref#1{algorithm~\ref{#1}}
% Reference to an algorithm, upper case.
\def\Algref#1{Algorithm~\ref{#1}}
\def\twoalgref#1#2{algorithms \ref{#1} and \ref{#2}}
\def\Twoalgref#1#2{Algorithms \ref{#1} and \ref{#2}}
% Reference to a part, lower case
\def\partref#1{part~\ref{#1}}
% Reference to a part, upper case
\def\Partref#1{Part~\ref{#1}}
\def\twopartref#1#2{parts \ref{#1} and \ref{#2}}

\def\ceil#1{\lceil #1 \rceil}
\def\floor#1{\lfloor #1 \rfloor}
\def\1{\bm{1}}
\newcommand{\train}{\mathcal{D}}
\newcommand{\valid}{\mathcal{D_{\mathrm{valid}}}}
\newcommand{\test}{\mathcal{D_{\mathrm{test}}}}

\def\eps{{\epsilon}}


% Random variables
\def\reta{{\textnormal{$\eta$}}}
\def\ra{{\textnormal{a}}}
\def\rb{{\textnormal{b}}}
\def\rc{{\textnormal{c}}}
\def\rd{{\textnormal{d}}}
\def\re{{\textnormal{e}}}
\def\rf{{\textnormal{f}}}
\def\rg{{\textnormal{g}}}
\def\rh{{\textnormal{h}}}
\def\ri{{\textnormal{i}}}
\def\rj{{\textnormal{j}}}
\def\rk{{\textnormal{k}}}
\def\rl{{\textnormal{l}}}
% rm is already a command, just don't name any random variables m
\def\rn{{\textnormal{n}}}
\def\ro{{\textnormal{o}}}
\def\rp{{\textnormal{p}}}
\def\rq{{\textnormal{q}}}
\def\rr{{\textnormal{r}}}
\def\rs{{\textnormal{s}}}
\def\rt{{\textnormal{t}}}
\def\ru{{\textnormal{u}}}
\def\rv{{\textnormal{v}}}
\def\rw{{\textnormal{w}}}
\def\rx{{\textnormal{x}}}
\def\ry{{\textnormal{y}}}
\def\rz{{\textnormal{z}}}

% Random vectors
\def\rvepsilon{{\mathbf{\epsilon}}}
\def\rvtheta{{\mathbf{\theta}}}
\def\rva{{\mathbf{a}}}
\def\rvb{{\mathbf{b}}}
\def\rvc{{\mathbf{c}}}
\def\rvd{{\mathbf{d}}}
\def\rve{{\mathbf{e}}}
\def\rvf{{\mathbf{f}}}
\def\rvg{{\mathbf{g}}}
\def\rvh{{\mathbf{h}}}
\def\rvu{{\mathbf{i}}}
\def\rvj{{\mathbf{j}}}
\def\rvk{{\mathbf{k}}}
\def\rvl{{\mathbf{l}}}
\def\rvm{{\mathbf{m}}}
\def\rvn{{\mathbf{n}}}
\def\rvo{{\mathbf{o}}}
\def\rvp{{\mathbf{p}}}
\def\rvq{{\mathbf{q}}}
\def\rvr{{\mathbf{r}}}
\def\rvs{{\mathbf{s}}}
\def\rvt{{\mathbf{t}}}
\def\rvu{{\mathbf{u}}}
\def\rvv{{\mathbf{v}}}
\def\rvw{{\mathbf{w}}}
\def\rvx{{\mathbf{x}}}
\def\rvy{{\mathbf{y}}}
\def\rvz{{\mathbf{z}}}

% Elements of random vectors
\def\erva{{\textnormal{a}}}
\def\ervb{{\textnormal{b}}}
\def\ervc{{\textnormal{c}}}
\def\ervd{{\textnormal{d}}}
\def\erve{{\textnormal{e}}}
\def\ervf{{\textnormal{f}}}
\def\ervg{{\textnormal{g}}}
\def\ervh{{\textnormal{h}}}
\def\ervi{{\textnormal{i}}}
\def\ervj{{\textnormal{j}}}
\def\ervk{{\textnormal{k}}}
\def\ervl{{\textnormal{l}}}
\def\ervm{{\textnormal{m}}}
\def\ervn{{\textnormal{n}}}
\def\ervo{{\textnormal{o}}}
\def\ervp{{\textnormal{p}}}
\def\ervq{{\textnormal{q}}}
\def\ervr{{\textnormal{r}}}
\def\ervs{{\textnormal{s}}}
\def\ervt{{\textnormal{t}}}
\def\ervu{{\textnormal{u}}}
\def\ervv{{\textnormal{v}}}
\def\ervw{{\textnormal{w}}}
\def\ervx{{\textnormal{x}}}
\def\ervy{{\textnormal{y}}}
\def\ervz{{\textnormal{z}}}

% Random matrices
\def\rmA{{\mathbf{A}}}
\def\rmB{{\mathbf{B}}}
\def\rmC{{\mathbf{C}}}
\def\rmD{{\mathbf{D}}}
\def\rmE{{\mathbf{E}}}
\def\rmF{{\mathbf{F}}}
\def\rmG{{\mathbf{G}}}
\def\rmH{{\mathbf{H}}}
\def\rmI{{\mathbf{I}}}
\def\rmJ{{\mathbf{J}}}
\def\rmK{{\mathbf{K}}}
\def\rmL{{\mathbf{L}}}
\def\rmM{{\mathbf{M}}}
\def\rmN{{\mathbf{N}}}
\def\rmO{{\mathbf{O}}}
\def\rmP{{\mathbf{P}}}
\def\rmQ{{\mathbf{Q}}}
\def\rmR{{\mathbf{R}}}
\def\rmS{{\mathbf{S}}}
\def\rmT{{\mathbf{T}}}
\def\rmU{{\mathbf{U}}}
\def\rmV{{\mathbf{V}}}
\def\rmW{{\mathbf{W}}}
\def\rmX{{\mathbf{X}}}
\def\rmY{{\mathbf{Y}}}
\def\rmZ{{\mathbf{Z}}}

% Elements of random matrices
\def\ermA{{\textnormal{A}}}
\def\ermB{{\textnormal{B}}}
\def\ermC{{\textnormal{C}}}
\def\ermD{{\textnormal{D}}}
\def\ermE{{\textnormal{E}}}
\def\ermF{{\textnormal{F}}}
\def\ermG{{\textnormal{G}}}
\def\ermH{{\textnormal{H}}}
\def\ermI{{\textnormal{I}}}
\def\ermJ{{\textnormal{J}}}
\def\ermK{{\textnormal{K}}}
\def\ermL{{\textnormal{L}}}
\def\ermM{{\textnormal{M}}}
\def\ermN{{\textnormal{N}}}
\def\ermO{{\textnormal{O}}}
\def\ermP{{\textnormal{P}}}
\def\ermQ{{\textnormal{Q}}}
\def\ermR{{\textnormal{R}}}
\def\ermS{{\textnormal{S}}}
\def\ermT{{\textnormal{T}}}
\def\ermU{{\textnormal{U}}}
\def\ermV{{\textnormal{V}}}
\def\ermW{{\textnormal{W}}}
\def\ermX{{\textnormal{X}}}
\def\ermY{{\textnormal{Y}}}
\def\ermZ{{\textnormal{Z}}}

% Vectors
\def\vzero{{\bm{0}}}
\def\vone{{\bm{1}}}
\def\vmu{{\bm{\mu}}}
\def\vtheta{{\bm{\theta}}}
\def\va{{\bm{a}}}
\def\vb{{\bm{b}}}
\def\vc{{\bm{c}}}
\def\vd{{\bm{d}}}
\def\ve{{\bm{e}}}
\def\vf{{\bm{f}}}
\def\vg{{\bm{g}}}
\def\vh{{\bm{h}}}
\def\vi{{\bm{i}}}
\def\vj{{\bm{j}}}
\def\vk{{\bm{k}}}
\def\vl{{\bm{l}}}
\def\vm{{\bm{m}}}
\def\vn{{\bm{n}}}
\def\vo{{\bm{o}}}
\def\vp{{\bm{p}}}
\def\vq{{\bm{q}}}
\def\vr{{\bm{r}}}
\def\vs{{\bm{s}}}
\def\vt{{\bm{t}}}
\def\vu{{\bm{u}}}
\def\vv{{\bm{v}}}
\def\vw{{\bm{w}}}
\def\vx{{\bm{x}}}
\def\vy{{\bm{y}}}
\def\vz{{\bm{z}}}

% Elements of vectors
\def\evalpha{{\alpha}}
\def\evbeta{{\beta}}
\def\evepsilon{{\epsilon}}
\def\evlambda{{\lambda}}
\def\evomega{{\omega}}
\def\evmu{{\mu}}
\def\evpsi{{\psi}}
\def\evsigma{{\sigma}}
\def\evtheta{{\theta}}
\def\eva{{a}}
\def\evb{{b}}
\def\evc{{c}}
\def\evd{{d}}
\def\eve{{e}}
\def\evf{{f}}
\def\evg{{g}}
\def\evh{{h}}
\def\evi{{i}}
\def\evj{{j}}
\def\evk{{k}}
\def\evl{{l}}
\def\evm{{m}}
\def\evn{{n}}
\def\evo{{o}}
\def\evp{{p}}
\def\evq{{q}}
\def\evr{{r}}
\def\evs{{s}}
\def\evt{{t}}
\def\evu{{u}}
\def\evv{{v}}
\def\evw{{w}}
\def\evx{{x}}
\def\evy{{y}}
\def\evz{{z}}

% Matrix
\def\mA{{\bm{A}}}
\def\mB{{\bm{B}}}
\def\mC{{\bm{C}}}
\def\mD{{\bm{D}}}
\def\mE{{\bm{E}}}
\def\mF{{\bm{F}}}
\def\mG{{\bm{G}}}
\def\mH{{\bm{H}}}
\def\mI{{\bm{I}}}
\def\mJ{{\bm{J}}}
\def\mK{{\bm{K}}}
\def\mL{{\bm{L}}}
\def\mM{{\bm{M}}}
\def\mN{{\bm{N}}}
\def\mO{{\bm{O}}}
\def\mP{{\bm{P}}}
\def\mQ{{\bm{Q}}}
\def\mR{{\bm{R}}}
\def\mS{{\bm{S}}}
\def\mT{{\bm{T}}}
\def\mU{{\bm{U}}}
\def\mV{{\bm{V}}}
\def\mW{{\bm{W}}}
\def\mX{{\bm{X}}}
\def\mY{{\bm{Y}}}
\def\mZ{{\bm{Z}}}
\def\mBeta{{\bm{\beta}}}
\def\mPhi{{\bm{\Phi}}}
\def\mLambda{{\bm{\Lambda}}}
\def\mSigma{{\bm{\Sigma}}}

% Tensor
\DeclareMathAlphabet{\mathsfit}{\encodingdefault}{\sfdefault}{m}{sl}
\SetMathAlphabet{\mathsfit}{bold}{\encodingdefault}{\sfdefault}{bx}{n}
\newcommand{\tens}[1]{\bm{\mathsfit{#1}}}
\def\tA{{\tens{A}}}
\def\tB{{\tens{B}}}
\def\tC{{\tens{C}}}
\def\tD{{\tens{D}}}
\def\tE{{\tens{E}}}
\def\tF{{\tens{F}}}
\def\tG{{\tens{G}}}
\def\tH{{\tens{H}}}
\def\tI{{\tens{I}}}
\def\tJ{{\tens{J}}}
\def\tK{{\tens{K}}}
\def\tL{{\tens{L}}}
\def\tM{{\tens{M}}}
\def\tN{{\tens{N}}}
\def\tO{{\tens{O}}}
\def\tP{{\tens{P}}}
\def\tQ{{\tens{Q}}}
\def\tR{{\tens{R}}}
\def\tS{{\tens{S}}}
\def\tT{{\tens{T}}}
\def\tU{{\tens{U}}}
\def\tV{{\tens{V}}}
\def\tW{{\tens{W}}}
\def\tX{{\tens{X}}}
\def\tY{{\tens{Y}}}
\def\tZ{{\tens{Z}}}


% Graph
\def\gA{{\mathcal{A}}}
\def\gB{{\mathcal{B}}}
\def\gC{{\mathcal{C}}}
\def\gD{{\mathcal{D}}}
\def\gE{{\mathcal{E}}}
\def\gF{{\mathcal{F}}}
\def\gG{{\mathcal{G}}}
\def\gH{{\mathcal{H}}}
\def\gI{{\mathcal{I}}}
\def\gJ{{\mathcal{J}}}
\def\gK{{\mathcal{K}}}
\def\gL{{\mathcal{L}}}
\def\gM{{\mathcal{M}}}
\def\gN{{\mathcal{N}}}
\def\gO{{\mathcal{O}}}
\def\gP{{\mathcal{P}}}
\def\gQ{{\mathcal{Q}}}
\def\gR{{\mathcal{R}}}
\def\gS{{\mathcal{S}}}
\def\gT{{\mathcal{T}}}
\def\gU{{\mathcal{U}}}
\def\gV{{\mathcal{V}}}
\def\gW{{\mathcal{W}}}
\def\gX{{\mathcal{X}}}
\def\gY{{\mathcal{Y}}}
\def\gZ{{\mathcal{Z}}}

% Sets
\def\sA{{\mathbb{A}}}
\def\sB{{\mathbb{B}}}
\def\sC{{\mathbb{C}}}
\def\sD{{\mathbb{D}}}
% Don't use a set called E, because this would be the same as our symbol
% for expectation.
\def\sF{{\mathbb{F}}}
\def\sG{{\mathbb{G}}}
\def\sH{{\mathbb{H}}}
\def\sI{{\mathbb{I}}}
\def\sJ{{\mathbb{J}}}
\def\sK{{\mathbb{K}}}
\def\sL{{\mathbb{L}}}
\def\sM{{\mathbb{M}}}
\def\sN{{\mathbb{N}}}
\def\sO{{\mathbb{O}}}
\def\sP{{\mathbb{P}}}
\def\sQ{{\mathbb{Q}}}
\def\sR{{\mathbb{R}}}
\def\sS{{\mathbb{S}}}
\def\sT{{\mathbb{T}}}
\def\sU{{\mathbb{U}}}
\def\sV{{\mathbb{V}}}
\def\sW{{\mathbb{W}}}
\def\sX{{\mathbb{X}}}
\def\sY{{\mathbb{Y}}}
\def\sZ{{\mathbb{Z}}}

% Entries of a matrix
\def\emLambda{{\Lambda}}
\def\emA{{A}}
\def\emB{{B}}
\def\emC{{C}}
\def\emD{{D}}
\def\emE{{E}}
\def\emF{{F}}
\def\emG{{G}}
\def\emH{{H}}
\def\emI{{I}}
\def\emJ{{J}}
\def\emK{{K}}
\def\emL{{L}}
\def\emM{{M}}
\def\emN{{N}}
\def\emO{{O}}
\def\emP{{P}}
\def\emQ{{Q}}
\def\emR{{R}}
\def\emS{{S}}
\def\emT{{T}}
\def\emU{{U}}
\def\emV{{V}}
\def\emW{{W}}
\def\emX{{X}}
\def\emY{{Y}}
\def\emZ{{Z}}
\def\emSigma{{\Sigma}}

% entries of a tensor
% Same font as tensor, without \bm wrapper
\newcommand{\etens}[1]{\mathsfit{#1}}
\def\etLambda{{\etens{\Lambda}}}
\def\etA{{\etens{A}}}
\def\etB{{\etens{B}}}
\def\etC{{\etens{C}}}
\def\etD{{\etens{D}}}
\def\etE{{\etens{E}}}
\def\etF{{\etens{F}}}
\def\etG{{\etens{G}}}
\def\etH{{\etens{H}}}
\def\etI{{\etens{I}}}
\def\etJ{{\etens{J}}}
\def\etK{{\etens{K}}}
\def\etL{{\etens{L}}}
\def\etM{{\etens{M}}}
\def\etN{{\etens{N}}}
\def\etO{{\etens{O}}}
\def\etP{{\etens{P}}}
\def\etQ{{\etens{Q}}}
\def\etR{{\etens{R}}}
\def\etS{{\etens{S}}}
\def\etT{{\etens{T}}}
\def\etU{{\etens{U}}}
\def\etV{{\etens{V}}}
\def\etW{{\etens{W}}}
\def\etX{{\etens{X}}}
\def\etY{{\etens{Y}}}
\def\etZ{{\etens{Z}}}

% The true underlying data generating distribution
\newcommand{\pdata}{p_{\rm{data}}}
% The empirical distribution defined by the training set
\newcommand{\ptrain}{\hat{p}_{\rm{data}}}
\newcommand{\Ptrain}{\hat{P}_{\rm{data}}}
% The model distribution
\newcommand{\pmodel}{p_{\rm{model}}}
\newcommand{\Pmodel}{P_{\rm{model}}}
\newcommand{\ptildemodel}{\tilde{p}_{\rm{model}}}
% Stochastic autoencoder distributions
\newcommand{\pencode}{p_{\rm{encoder}}}
\newcommand{\pdecode}{p_{\rm{decoder}}}
\newcommand{\precons}{p_{\rm{reconstruct}}}

\newcommand{\laplace}{\mathrm{Laplace}} % Laplace distribution

\newcommand{\E}{\mathbb{E}}
\newcommand{\Ls}{\mathcal{L}}
\newcommand{\R}{\mathbb{R}}
\newcommand{\emp}{\tilde{p}}
\newcommand{\lr}{\alpha}
\newcommand{\reg}{\lambda}
\newcommand{\rect}{\mathrm{rectifier}}
\newcommand{\softmax}{\mathrm{softmax}}
\newcommand{\sigmoid}{\sigma}
\newcommand{\softplus}{\zeta}
\newcommand{\KL}{D_{\mathrm{KL}}}
\newcommand{\Var}{\mathrm{Var}}
\newcommand{\standarderror}{\mathrm{SE}}
\newcommand{\Cov}{\mathrm{Cov}}
% Wolfram Mathworld says $L^2$ is for function spaces and $\ell^2$ is for vectors
% But then they seem to use $L^2$ for vectors throughout the site, and so does
% wikipedia.
\newcommand{\normlzero}{L^0}
\newcommand{\normlone}{L^1}
\newcommand{\normltwo}{L^2}
\newcommand{\normlp}{L^p}
\newcommand{\normmax}{L^\infty}

\newcommand{\parents}{Pa} % See usage in notation.tex. Chosen to match Daphne's book.

\DeclareMathOperator*{\argmax}{arg\,max}
\DeclareMathOperator*{\argmin}{arg\,min}

\DeclareMathOperator{\sign}{sign}
\DeclareMathOperator{\Tr}{Tr}
\let\ab\allowbreak

%%%%% NEW MATH DEFINITIONS %%%%%

\usepackage{amsmath,amsfonts,bm}

% Mark sections of captions for referring to divisions of figures


% Figure reference, lower-case.
\def\figref#1{figure~\ref{#1}}
% Figure reference, capital. For start of sentence
\def\Figref#1{Figure~\ref{#1}}
\def\twofigref#1#2{figures \ref{#1} and \ref{#2}}
\def\quadfigref#1#2#3#4{figures \ref{#1}, \ref{#2}, \ref{#3} and \ref{#4}}
% Section reference, lower-case.
\def\secref#1{section~\ref{#1}}
% Section reference, capital.
\def\Secref#1{Section~\ref{#1}}
% Reference to two sections.
\def\twosecrefs#1#2{sections \ref{#1} and \ref{#2}}
% Reference to three sections.
\def\secrefs#1#2#3{sections \ref{#1}, \ref{#2} and \ref{#3}}
% Reference to an equation, lower-case.
\def\eqref#1{equation~\ref{#1}}
% Reference to an equation, upper case
\def\Eqref#1{Equation~\ref{#1}}
% A raw reference to an equation---avoid using if possible
\def\plaineqref#1{\ref{#1}}
% Reference to a chapter, lower-case.
\def\chapref#1{chapter~\ref{#1}}
% Reference to an equation, upper case.
\def\Chapref#1{Chapter~\ref{#1}}
% Reference to a range of chapters
\def\rangechapref#1#2{chapters\ref{#1}--\ref{#2}}
% Reference to an algorithm, lower-case.
\def\algref#1{algorithm~\ref{#1}}
% Reference to an algorithm, upper case.
\def\Algref#1{Algorithm~\ref{#1}}
\def\twoalgref#1#2{algorithms \ref{#1} and \ref{#2}}
\def\Twoalgref#1#2{Algorithms \ref{#1} and \ref{#2}}
% Reference to a part, lower case
\def\partref#1{part~\ref{#1}}
% Reference to a part, upper case
\def\Partref#1{Part~\ref{#1}}
\def\twopartref#1#2{parts \ref{#1} and \ref{#2}}

\def\ceil#1{\lceil #1 \rceil}
\def\floor#1{\lfloor #1 \rfloor}
\def\1{\bm{1}}


\def\eps{{\epsilon}}


\let\ab\allowbreak


\definecolor{myred}{RGB}{255, 67, 20}
\definecolor{myblue}{RGB}{51, 51, 255}
\definecolor{mygreen}{RGB}{0, 128, 0}

\begin{document}

\title{SNN 本地学习的逐层信息保留与损失定位（2026-10）}

\author{}
\date{}

\maketitle

\begin{abstract}
本地学习算法（OTTT、NDOT、S-TLLR、TESS）通过时间局部化与空间局部化降低训练开销，但它们各自优化局部目标，没有像 BPTT 那样统一的全局梯度信号。因此在现有评估中，只能看到最终准确率，无法判断信息在哪一层损失、损失了多少。本方案提出在层与层之间引入中间量，逐层量化本地学习的信息保留程度与训练效果。具体包括三个研究方向：逐层探针、逐层表示相似度、逐层梯度对齐。三者分别从"最终效用、学习结果、学习过程"三个层面定位信息损失，共同回答"信息在哪一层损失、损失多少、为什么损失"。

本文是 2026-10 月度版本，在沿用 2026-09 核心问题与三个研究方向的基础上，补充：已完成工作、分阶段下一步计划、以及组会报告汇总。
\end{abstract}

\section*{文章汇总}

RTRL: \textbf{A learning algorithm for continually running fully recurrent neural networks} \cite{williams1989learning};

E-prop: \textbf{A solution to the learning dilemma for recurrent networks of spiking neurons} \cite{bellec_solution_2020};

OTTT: \textbf{Online Training Through Time for Spiking Neural Networks} \cite{xiao_online_2022};

NDOT: \textbf{NDOT: Neuronal Dynamics-based Online Training for Spiking Neural Networks} \cite{jiang_ndot_nodate};

S-TLLR: \textbf{S-TLLR: STDP-inspired Temporal Local Learning Rule for Spiking Neural Networks} \cite{apolinario_s-tllr_2024};

TESS: \textbf{TESS: A Scalable Temporally and Spatially Local Learning Rule for Spiking Neural Networks} \cite{apolinario_tess_2025}.

\section{研究背景与核心问题}

\subsection{本地学习缺少全局梯度}

本地学习算法是降低 SNN 训练开销的主要途径，其核心思路是在时间维度和空间维度上切断 BPTT 的全局反向链，用局部可得的信息近似梯度。四种代表性算法在“丢什么、补什么”上依次递进：

OTTT \cite{xiao_online_2022} 在时间维度上丢弃脉冲介导路径（spike-mediated path），即时间 Jacobian 中经过脉冲 $s_t$ 的 $-\lambda V_{th}D_\Psi$ 项，等价于复位路径不反传；只保留固定衰减 $\lambda$，把时间依赖近似为 $A_t\approx\lambda I$。

NDOT \cite{jiang_ndot_nodate} 保留时间依赖，但用神经元动力学系数 $e^l[t]=\frac{u^l[t]-V_{th}s^l[t]}{u^l[t-1]-V_{th}s^l[t-1]}$ 替代 OTTT 的固定 $\lambda$，使时间反传系数随膜电位动态变化，比 OTTT 更接近完整 BPTT。

S-TLLR \cite{apolinario_s-tllr_2024} 转向 STDP 风格的时间信用分配，用因果（pre-before-post）与非因果（post-before-pre）两条资格迹的瞬时组合替代显式时间反传，把时间局部性建立在突触前/后脉冲的时序关系上，而非膜电位递推。

TESS \cite{apolinario_tess_2025} 在 S-TLLR 的时间局部性基础上，进一步用本地生成的学习信号 $m^{(l)}[t]=B^{(l)\top}(f(B^{(l)}o^{(l)}[t])-y^*)$ 替代跨层误差反传，同时实现时间与空间全局部。

四者的共同点是：每层（或每个突触）只依据局部可得的信息优化一个局部目标，网络内部不存在像 BPTT 那样贯穿全网的统一梯度。因此在现有评估框架下，只能看到最终准确率，无法回答信息在哪一层损失、损失多少。

BPTT 的梯度有明确的全局结构：误差从输出层出发，沿空间（层到层）和时间（$t$ 到 $t-1$）两个方向逐层、逐时刻地传播，最终落到每个权重上。因此对 BPTT 而言，"信息在哪一层损失"至少在梯度层面是可定义的——可以逐层检查反向传播到该层的误差信号大小。而本地学习切断了这一全局链条：每层各自优化局部目标，层与层之间没有统一的误差信号可追踪，于是"信息在哪一层损失、损失多少"这一问题在现有评估框架下没有答案。

\subsection{核心研究问题}

目前的评估方式停留在"看最终准确率"：给每个本地学习算法在某个任务上报告一个准确率，谁高谁好。但这无法回答三个关键问题：

\begin{enumerate}[leftmargin=1.5em,itemsep=0pt]
\item 信息在哪一层损失：本地学习相对 BPTT 的性能差距，究竟来自浅层还是深层？
\item 损失了多少：每一层各自保留了多少任务信息、丢弃了多少？
\item 为什么损失：损失来自时间近似的结构性误差、空间目标替换，还是两者耦合？
\end{enumerate}

核心研究问题：\textbf{在层与层之间引入中间量，逐层量化本地学习的信息保留程度与训练效果，从而定位"信息在哪一层损失、损失多少"及其来源。}

\section{三个研究方向}

三个方向分别从不同层面切入同一个核心问题，构成从"因"到"果"的链条：梯度对齐度量学习过程本身是否偏离，表示相似度度量学习结果是否被扭曲，逐层探针度量最终任务信息是否保留。下面逐一说明每个方向要回答什么问题、为什么需要、如何定义指标，以及与另外两个方向的关系。

\subsection{方向一：逐层探针}

\textbf{要回答的问题}：第 $l$ 层的表示本身携带了多少关于任务标签 $y$ 的信息？即"任务信息在到达第 $l$ 层时还保留多少"。

\textbf{为什么需要}：这是三个方向里最直接的"信息损失"度量。如果第 $l$ 层的表示已经能够以高精度解码出任务标签，说明信息在到达该层时尚未丢失；如果解码精度很低，说明信息在该层之前（或该层内部）已经丢失。逐层扫描即可定位损失发生的位置。

\textbf{探针的挂法}：在每一层输出之后挂一个独立的、小的解码器（探针），只以该层的表示作为输入去预测任务标签。主网络的权重在挂探针时保持冻结，探针只在冻结的表示之上训练。具体地，第 $l$ 层在 $T$ 个时间步的脉冲序列 $s^l[1],\dots,s^l[T]$ 先聚合为一个向量表示 $z^l$，例如取时间平均（firing rate）
\begin{equation}
z^l = \frac{1}{T}\sum_{t=1}^{T} s^l[t]\in\mathbb R^{n^l}.
\end{equation}
探针取线性分类器
\begin{equation}
f^l_{\mathrm{probe}}(z^l) = \operatorname{softmax}\big(P^l z^l + c^l\big),
\qquad
P^l\in\mathbb R^{C\times n^l},\ c^l\in\mathbb R^C,
\end{equation}
其中 $C$ 是类别数。探针参数 $P^l, c^l$ 通过最小化交叉熵（或其他任务损失）在训练集上拟合，训练时不回传梯度到主网络。

\textbf{探针准确率的计算}：在验证集上，用训练好的探针对每个样本 $n$ 做预测 $\hat y_n = \arg\max_c f^l_{\mathrm{probe}}(z^l_n)$，探针准确率定义为
\begin{equation}
\operatorname{Acc}^l_{\mathrm{probe}} = \frac{1}{N}\sum_{n=1}^{N}\mathbb I[\hat y_n = y_n].
\end{equation}

\textbf{归一化信息保留率}：为了跨层比较，把准确率归一化到"从随机猜测到输出层"的区间，定义
\begin{equation}
\rho^l = \frac{\operatorname{Acc}^l_{\mathrm{probe}} - \operatorname{Acc}_{\mathrm{chance}}}{\operatorname{Acc}^L_{\mathrm{probe}} - \operatorname{Acc}_{\mathrm{chance}}},
\end{equation}
其中 $\operatorname{Acc}_{\mathrm{chance}}=1/C$ 是随机猜测水平，$\operatorname{Acc}^L_{\mathrm{probe}}$ 是输出层（第 $L$ 层）探针的准确率，作为"该任务可达到的信息上限"。

\textbf{如何用曲线判断信息在层间的分布}：画出 $\rho^l$ 对层索引 $l$ 的曲线。$\rho^l$ 越接近 $1$，说明该层表示保留的任务信息越接近上限；曲线中下降最快的区间，就是信息损失最集中的层。将本地学习算法与 BPTT 的两条 $\rho^l$ 曲线画在同一张图上：两者差距最大的层，就是该本地学习算法相对 BPTT 信息损失最大的层。

\textbf{与另外两个方向的关系}：探针回答"信息还在不在"，属于"最终效用"层面，但它只看结果、不解释原因——它无法说明"为什么这一层丢了信息"。要回答"为什么"，需要表示相似度（看表示结构是否被扭曲）和梯度对齐（看优化方向是否偏离）来补充。

\subsection{方向二：逐层表示相似度}

\textbf{要回答的问题}：本地学习训练出的第 $l$ 层表示，与 BPTT 训练出的第 $l$ 层表示，在几何结构上有多像？

\textbf{为什么需要}：探针只能说明"信息是否还在"，但不能说明"表示本身是否被扭曲"。两个不同的表示可能保留同样的任务信息（探针准确率相同），但几何结构完全不同——例如两者相差一个复杂的非线性变换。这种结构差异虽然不降低当前任务的解码精度，却会影响泛化、迁移和与后续层的配合。因此需要直接比较表示本身，而不只是比较它们能解码出什么。

\textbf{表示相似度指标的定义}：对第 $l$ 层，收集 $N$ 个样本的表示，分别排成矩阵：本地学习给出 $Z^l_{\mathrm{local}}\in\mathbb R^{N\times n^l}$，BPTT 给出 $Z^l_{\mathrm{BPTT}}\in\mathbb R^{N\times n^l}$，其中第 $n$ 行是第 $n$ 个样本的表示 $z^l_n$。比较两个表示，本质上是比较它们各自定义的"样本间几何结构"。用 Gram 矩阵刻画这一结构：$K^l_{\mathrm{local}} = Z^l_{\mathrm{local}}(Z^l_{\mathrm{local}})^\top$ 的第 $(i,j)$ 元是样本 $i$ 与样本 $j$ 表示的内积，度量两样本在该表示下的相似程度；$K^l_{\mathrm{BPTT}}$ 同理。为了让比较不受表示均值（全局平移）的影响，先对 Gram 矩阵做中心化：记 $H = I - \frac{1}{N}\mathbf{1}\mathbf{1}^{\top}$（$\mathbf{1}$ 为全 1 向量，$I$ 为单位阵），定义
\begin{equation}
\tilde K^l_{\mathrm{local}} = H K^l_{\mathrm{local}} H,
\qquad
\tilde K^l_{\mathrm{BPTT}} = H K^l_{\mathrm{BPTT}} H.
\end{equation}
然后定义第 $l$ 层的表示相似度为两个中心化 Gram 矩阵的（归一化）内积：
\begin{equation}
\operatorname{sim}^l = \frac{\langle \tilde K^l_{\mathrm{local}}, \tilde K^l_{\mathrm{BPTT}}\rangle_F}{\|\tilde K^l_{\mathrm{local}}\|_F\,\|\tilde K^l_{\mathrm{BPTT}}\|_F}
= \frac{\operatorname{tr}\big(\tilde K^l_{\mathrm{local}}\,\tilde K^l_{\mathrm{BPTT}}\big)}{\|\tilde K^l_{\mathrm{local}}\|_F\,\|\tilde K^l_{\mathrm{BPTT}}\|_F},
\end{equation}
其中 $\langle A,B\rangle_F=\operatorname{tr}(A^\top B)$ 是 Frobenius 内积，$\|\cdot\|_F$ 是 Frobenius 范数。

\textbf{取值范围与含义}：$\operatorname{sim}^l$ 本质上是两个矩阵（视为高维向量）之间夹角的余弦，取值范围为 $[-1,1]$，实际中通常落在 $[0,1]$。$\operatorname{sim}^l=1$ 表示两种表示的"样本间几何结构"完全一致（只可能相差一个保持该结构的正交变换）；$\operatorname{sim}^l$ 越低，表示本地学习的表示相对 BPTT 偏离越大。用 Gram 矩阵而非原始表示的原因在于：它只关心"样本之间的关系结构"，对表示自身的旋转、以及（经过归一化后）缩放不敏感，从而把比较聚焦在"结构是否被扭曲"上。

\textbf{逐层曲线的读法}：画出 $\operatorname{sim}^l$ 对 $l$ 的曲线。浅层通常相似度较高（输入编码尚未被明显改变），越往深层差异可能越大。相似度快速下降的层，就是本地学习相对 BPTT 表示被"扭曲"最严重的层。

\textbf{与另外两个方向的关系}：表示相似度回答"表示结构变没变"，属于"学习结果"层面，是探针（最终效用）与梯度对齐（学习过程）之间的中间环节。它比探针更能定位"在哪一层几何结构开始偏离"，但相似度下降并不必然等于任务信息丢失——有时只是换了一个等价表示。因此需要与探针交叉验证：相似度低且探针准确率也低，才是真正"有害"的信息损失。

\subsection{方向三：逐层梯度对齐}

\textbf{要回答的问题}：本地学习在第 $l$ 层产生的梯度（或等效更新方向）与 BPTT 在第 $l$ 层的梯度，方向是否一致？

\textbf{为什么需要}：探针和表示相似度都是"事后"看结果，梯度对齐看的是"学习过程"本身。即使某层当前表示与 BPTT 相似，如果本地学习的梯度方向错误，那么在后续训练中该层会逐渐偏离。梯度对齐直接度量"本地学习是否正朝着 BPTT 会走的方向优化"，因此是三个方向里最接近"原因"的度量。

\textbf{梯度对齐的定义}：把第 $l$ 层的权重 $W^{l\leftarrow l-1}$ 展平为一个向量，本地学习给出的梯度记为 $g^l_{\mathrm{local}}\in\mathbb R^d$，BPTT 给出的梯度记为 $g^l_{\mathrm{BPTT}}\in\mathbb R^d$。二者的方向一致性用余弦相似度度量：
\begin{equation}
\cos^l = \frac{\langle g^l_{\mathrm{local}}, g^l_{\mathrm{BPTT}}\rangle}{\|g^l_{\mathrm{local}}\|\,\|g^l_{\mathrm{BPTT}}\|}
\in[-1,1].
\end{equation}
$\cos^l=1$ 表示方向完全一致，$=0$ 表示正交，$<0$ 表示方向相反。

\textbf{平行/垂直分量分解}：为了把"方向一致"进一步拆成"有效的平行部分"与"有害的垂直部分"，把 $g^l_{\mathrm{local}}$ 沿 $g^l_{\mathrm{BPTT}}$ 方向做正交分解：
\begin{equation}
g^l_{\parallel} = \frac{\langle g^l_{\mathrm{local}}, g^l_{\mathrm{BPTT}}\rangle}{\|g^l_{\mathrm{BPTT}}\|^2}\,g^l_{\mathrm{BPTT}},
\qquad
g^l_{\perp} = g^l_{\mathrm{local}} - g^l_{\parallel}.
\end{equation}
$g^l_{\parallel}$ 是本地梯度中真正沿着 BPTT 方向、能推动正确优化的分量；$g^l_{\perp}$ 是垂直于 BPTT 方向、与正确优化无关（甚至有害）的分量。相应地定义两个标量：
\begin{equation}
c^l_{\parallel} = \frac{\langle g^l_{\mathrm{local}}, g^l_{\mathrm{BPTT}}\rangle}{\|g^l_{\mathrm{BPTT}}\|^2},
\qquad
R^l_{\perp} = \frac{\|g^l_{\perp}\|}{\|g^l_{\mathrm{BPTT}}\|}.
\end{equation}
$c^l_{\parallel}$ 是"本地梯度中有效的平行部分占 BPTT 梯度的比例"，$R^l_{\perp}$ 是"垂直偏差相对 BPTT 梯度的大小"。

\textbf{逐层曲线的读法}：画出 $\cos^l$（或 $c^l_{\parallel}$、$R^l_{\perp}$）对 $l$ 的曲线。$\cos^l$ 低、或 $R^l_{\perp}$ 高的层，就是本地近似在该层引入的梯度误差最大的层，也就是最可能"带偏"该层优化的位置。

\textbf{与另外两个方向的关系}：梯度对齐回答"优化方向对不对"，属于"学习过程"层面，是三者的"因"。梯度方向错了会导致表示偏离（方向二），进而导致任务信息丢失（方向一）。三者构成"梯度（过程）$\to$ 表示（结果）$\to$ 任务信息（效用）"的因果链，共同回答核心问题。

\subsection{三个方向的关系与整体定位}

三个方向构成一条从因到果的链条，共同回答"信息在哪一层损失、损失多少、为什么损失"：

\begin{enumerate}[leftmargin=1.5em,itemsep=0pt]
\item \textbf{梯度对齐}（学习过程）：定位"哪一层的优化方向被本地近似带偏"；
\item \textbf{表示相似度}（学习结果）：定位"哪一层的表示结构被本地学习扭曲"；
\item \textbf{逐层探针}（最终效用）：定位"哪一层的任务信息真正丢失"。
\end{enumerate}

定位流程是：先用探针找到信息在哪一层损失（结果层），再用表示相似度确认表示在哪一层被扭曲（结构层），最后用梯度对齐解释为什么在这一层损失（原因层）。三者互补——探针只看结果、梯度只看过程，表示相似度居中承上启下——单独使用任何一个都不足以完整回答核心问题。

\section{已完成工作（2026-09 $\to$ 2026-10）}

本节记录从 2026-09 到 2026-10 之间已经完成的工作，它们是下一步计划的起点。

\subsection{研究计划文档}

确立了核心研究问题与三个研究方向（见上文），形成一份可长期迭代的研究计划。核心问题是"本地学习每层优化局部目标、无统一全局梯度，无法定位信息在哪一层损失"，三个方向（逐层探针 / 逐层表示相似度 / 逐层梯度对齐）分别从"最终效用 / 学习结果 / 学习过程"三个层面切入。

\subsection{Project 基础代码}

在 \texttt{Project} 目录下搭起了最小可运行基础：

\begin{itemize}[leftmargin=1.5em,itemsep=2pt]
\item LIF 神经元（前向 Heaviside + 反向 surrogate）、多层 SNN（\texttt{forward} / \texttt{forward\_states} / \texttt{param\_vector}）；
\item surrogate 接 autograd：自定义 \texttt{SpikeFunction}（前向 Heaviside、反向 fast-sigmoid $\Psi'$），已接入 \texttt{LIF.step}，单样本前向 + 反向贯通；
\item 数据与脉冲编码（泊松编码）、统一训练循环 \texttt{train\_epoch()}、指标接口 \texttt{evaluate()}；
\item \texttt{forward\_check.py}（前向手算 + surrogate 反向贯通）、\texttt{stage0\_check.py}（阶段 0 全链路），均跑通（PASS）；
\item 尚未实现：真实 DVS 数据加载（留 TODO）。
\end{itemize}

\subsection{PPT 笔记推导}

重写了 \texttt{SNNTrainingSlideNotes.tex} 中 BPTT、OTTT、NDOT 三个 section：

\begin{itemize}[leftmargin=1.5em,itemsep=2pt]
\item 统一为 Step 1--8 的详细推导，每步含完整公式 + 解释 + 符号来源 + 推导方向；
\item 明确三者关系链：BPTT 给出完整时间依赖 $\epsilon^l[i]=\lambda I-\lambda V_{th}D_{\Psi,i}$，OTTT 丢弃复位项用固定 $\lambda I$ 近似，NDOT 用动态 $e^l[t]$ 替代固定 $\lambda$。
\end{itemize}

\subsection{冻结的关键约定}

在读论文与推导过程中冻结了若干关键约定（详见《SNN 在线本地学习：细节问题清单》）：

\begin{itemize}[leftmargin=1.5em,itemsep=2pt]
\item LIF 复位形式：统一用形式 A，$u_{t+1}=\lambda(u_t-V_{th}s_t)+I$；
\item 层索引：统一用 $W^{l\leftarrow l-1}$；
\item 时间索引：$\epsilon^l[i]$ 内的 spike Jacobian 用 $D_{\Psi,i}$；
\item 资格迹粒度：e-prop 每突触（$O(Ln^2)$），OTTT/NDOT/S-TLLR/TESS 每层（$O(Ln)$）。
\end{itemize}

\subsection{Project 文件手册与函数注释}

在 \texttt{Project} 目录下新增 \texttt{MANUAL.md}，记录项目现状盘点、逐文件与逐函数说明、数学符号与代码对应表；同时为 Project 下所有 \texttt{.py} 文件的函数/类补齐了 docstring 与关键行内注释（不改逻辑、不改签名）。

\subsection{阶段 1：BPTT-SG（已完成）}

手写 BPTT-SG 前向/反向并跑通收敛曲线（\texttt{python -m sanity.gradient\_check} 全 PASS）：

\begin{itemize}[leftmargin=1.5em,itemsep=2pt]
\item \textbf{手写 BPTT-SG}（\texttt{algorithm/bptt\_sg.py}）：\texttt{bptt\_sg\_forward}（纯数值前向返回每层 $u,s$）、\texttt{bptt\_sg\_backward}（手写反向递推返回每层 $dW/db$）、\texttt{bptt\_sg\_step}（前向 + 损失 + 反向）。支持 detached-reset（默认，$A_t=\lambda I$）与 full（$A_t=\lambda I-\lambda V_{th}D_\Psi$）两种模式；
\item \textbf{Gate A/B/C 梯度验证}（\texttt{sanity/gradient\_check.py}）：Gate A 验时间 Jacobian、Gate B 验空间 Jacobian、Gate C 验最终梯度，手写 vs autograd 相对误差 $<10^{-5}$（实测 $10^{-7}\sim10^{-8}$），detached 与 full 各一遍；
\item \textbf{收敛曲线}：full BPTT 在 256 样本合成数据上 last3 均值 $>0.85$；detached-reset 作对照；
\item \textbf{符号对照备忘录}：新增 \texttt{SYMBOL\_MAPPING.md}，统一论文记法（$W^{l\leftarrow l-1}$，1-based）与代码记法（\texttt{weights[l]}，0-based）差 1 层的映射，修正 \texttt{MANUAL.md} 符号表偏移；
\item \textbf{接口扩展}：\texttt{LIF.step} / \texttt{MLPSNN.forward} / \texttt{forward\_states} 加 \texttt{detach\_reset} 参数（默认 False，向后兼容）。
\end{itemize}

\section{下一步计划}

按阶段推进，每阶段有明确目标与任务。核心原则：先搭架子，架子上只挂一个能跑穿的东西（BPTT-SG），再逐个挂载其他算法；探针、CKA、梯度对齐的钩子必须在阶段 2 焊上。

\subsection{阶段 0：搭建实验骨架}

\textbf{目标}：搭一套通用实验框架，架子上只挂一个能跑到头（前向 + 反向贯通）的东西。（本阶段已基本完成，见"已完成工作"；仅真实 DVS 数据加载留 TODO）

\textbf{任务}：
\begin{enumerate}[leftmargin=1.5em,itemsep=0pt]
\item 数据与脉冲编码模块（DVS Gesture 为主，CIFAR10-DVS 为辅，$T=10$）；
\item 统一网络骨架 \texttt{SNN(nn.Module)}，支持返回每层 $u, s$ 与固定参数顺序；
\item 统一训练循环 \texttt{train\_epoch()}，支持 Accumulated / Online 两种模式；
\item 指标接口 \texttt{evaluate()}，预留探针、CKA、梯度对齐的调用点；
\item surrogate 接 autograd：自定义 \texttt{autograd.Function}（前向 Heaviside、反向 $\Psi'$），让单样本前向 + 反向贯通。（已完成）
\end{enumerate}

\subsection{阶段 1：挂 BPTT-SG，跑穿上限（已完成）}

\textbf{目标}：把 BPTT-SG 作为唯一算法插入骨架，跑出一条完整收敛曲线。

\textbf{状态}：已达成。手写 BPTT-SG 前向/反向 + Gate A/B/C 梯度验证（相对误差 $<10^{-5}$）+ 收敛曲线（合成数据 last3 均值 $>0.85$）全部 PASS（详见「已完成工作 $\to$ 阶段 1」）。

\textbf{任务}：
\begin{enumerate}[leftmargin=1.5em,itemsep=0pt]
\item 通过 gradcheck（Gate A/B/C）验证手写 BPTT-SG 梯度与 autograd 一致（相对误差 $<10^{-5}$）—— 首任务，先把梯度算对；（已完成）
\item 手写 BPTT-SG 前向与反向，用形式 A，支持 detached-reset 与 full 两种模式（复用阶段 0 已接好的 surrogate）；（已完成）
\item 跑通 DVS Gesture，记录每层 $u, s$ 与梯度。（留待阶段 5 前接入真实数据）
\end{enumerate}

\subsection{阶段 2：焊上探针、CKA、梯度对齐的钩子}

\textbf{目标}：在写其他算法前，先把测量工具装好。

\textbf{任务}：
\begin{enumerate}[leftmargin=1.5em,itemsep=0pt]
\item 逐层探针模块：每层挂独立线性探针，不参与主网络训练；
\item 逐层 CKA 模块：收集每层表示矩阵，计算与 BPTT 的 CKA；
\item 逐层梯度对齐模块：收集每层梯度，计算与 BPTT 的余弦相似度；
\item 在 BPTT 自身上验证工具本身：$\mathrm{CKA}=1$，$\cos=1$。
\end{enumerate}

\subsection{阶段 3：挂载 OTTT / NDOT / e-prop}

\textbf{目标}：同一框架下逐个挂载其他算法。

\textbf{任务}：
\begin{enumerate}[leftmargin=1.5em,itemsep=0pt]
\item OTTT：实现 $\hat{a}^{l-1}[t]$ 前向递推与瞬时梯度，验证递推与闭式一致；
\item NDOT：实现动态系数 $e^l[t]$ 与动态迹递推，注意 $t=1$ 特例；
\item e-prop：实现 ideal factorization 与 eligibility 局部递推，验证 $\frac{dE}{dW_{ji}}=\sum_t L_j^t e_{ji}^t$。
\end{enumerate}

\subsection{阶段 4：接入 S-TLLR / TESS}

\textbf{目标}：外部验证，优先使用作者公开代码。

\textbf{任务}：
\begin{enumerate}[leftmargin=1.5em,itemsep=0pt]
\item 下载 S-TLLR 与 TESS 作者代码，在统一数据下复现论文报告的数字；
\item 逐层测量三个指标，确保测量口径一致。
\end{enumerate}

\subsection{阶段 5：逐层评估实验}

\textbf{目标}：回答核心研究问题。

\textbf{任务}：
\begin{enumerate}[leftmargin=1.5em,itemsep=0pt]
\item 逐层画三张图：探针准确率、CKA、梯度对齐；
\item 归因分析：temporal 近似（OTTT vs NDOT）、spatial 局部化（S-TLLR vs TESS）、目标替换（BPTT vs TESS）；
\item Pareto 分析：Memory--Fidelity、Time--Fidelity、Communication--Fidelity 三张散点图。
\end{enumerate}

\subsection{阶段 6：与硬件方向对接}

\textbf{目标}：把逐层信息分布整理成硬件可参考的规格。

\textbf{任务}：
\begin{enumerate}[leftmargin=1.5em,itemsep=0pt]
\item 状态清单：每个算法需存储哪些状态、状态生命周期；
\item 通信清单：每层通信需求、哪些可本地化；
\item 更新清单：更新频率、更新模式、权重写入次数；
\item 与硬件同学讨论，定期反馈可行性。
\end{enumerate}

\subsection{长期方向：基于逐层定位的算法改进}

\textbf{定位}：本节不是当前研究计划的阶段（不编入阶段 0--6），而是阶段 6 之后、依赖逐层评估结论才成立的长期方向，作为后续文章的种子。

\textbf{前提}：阶段 2 焊好探针 / CKA / 梯度对齐三个测量工具、阶段 3--5 完成各本地学习算法的逐层评估之后，才有依据开展。

\textbf{落点}：用三个指标共同定位「某层的主要损失来源」——逐层探针找信息在哪一层丢失、逐层 CKA 看该层表示相对 BPTT 是否被扭曲、逐层梯度对齐看该层优化方向是否偏离；一旦确定损失集中在第 $l$ 层、且主要来自时间近似或空间局部化，就针对该层改进 OTTT / NDOT 的局部近似（例如修正该层的资格迹递推或学习信号），而非空泛地提出新算法。

\section{组会报告汇总}

本节按日期记录每次组会，格式统一为"简述 + 关键结论 + 待讨论问题"。

\subsection{2026-10-08 组会}

\textbf{汇报主题}：SNN 本地学习的逐层信息保留与损失定位——研究计划与阶段进展。

本节为本次组会的完整讲稿，逐页给出「页标题 / 幻灯片内容 / 口播要点」，可直接照着做 PPT 与口头汇报。本次覆盖：研究背景与核心问题、三个研究方向、阶段 0--1 已完成工作、下一步计划，以及一个后续方向。

\textbf{第 1 页 · 封面}

\begin{itemize}[leftmargin=1.5em,itemsep=0pt]
\item \textbf{幻灯片}：题目「SNN 本地学习的逐层信息保留与损失定位」，副标题「研究计划与阶段进展」，汇报人、日期（2026-10-08）。
\item \textbf{口播}：主题：SNN 本地学习的逐层信息保留与损失定位；内容：研究计划 + 阶段 0/1 已完成进展。
\end{itemize}

\textbf{第 2 页 · 核心问题：本地学习看不到「哪一层损失」}

\begin{itemize}[leftmargin=1.5em,itemsep=0pt]
\item \textbf{幻灯片}：左侧列 OTTT / NDOT / S-TLLR / TESS 四个本地学习算法，右侧画 BPTT 贯穿全网的全局梯度链，中间用「断裂」示意本地学习切断了这条链；底部一句话核心问题。
\item \textbf{口播}：四个本地学习算法：每层只按局部信息优化局部目标，无全局梯度链；\textit{后果}：现有评估只能报最终准确率，无法回答在哪层损失、损失多少、为什么损失；本研究切入点：在层间引入中间量，逐层量化。
\end{itemize}

\textbf{第 3 页 · 三个研究方向：三把尺子}

\begin{itemize}[leftmargin=1.5em,itemsep=0pt]
\item \textbf{幻灯片}：三块并列——逐层探针（量任务信息还在不在）、逐层表示相似度 CKA（量表示几何有没有被扭曲）、逐层梯度对齐（量优化方向偏没偏）；下方一条因果链「梯度（过程）$\to$ 表示（结果）$\to$ 任务信息（效用）」。
\item \textbf{口播}：逐层探针：量任务信息还保留多少（效用）；逐层 CKA：量表示几何相对 BPTT 是否被扭曲（结果）；逐层梯度对齐：量优化方向是否偏离（过程）；三者构成「梯度 $\to$ 表示 $\to$ 任务信息」因果链。
\end{itemize}

\textbf{第 4 页 · 整体路线图（阶段 0--6）}

\begin{itemize}[leftmargin=1.5em,itemsep=0pt]
\item \textbf{幻灯片}：横向流程「阶段 0 搭骨架 $\to$ 阶段 1 跑穿 BPTT-SG $\to$ 阶段 2 焊测量钩子 $\to$ 阶段 3 挂 OTTT/NDOT/e-prop $\to$ 阶段 4 接 S-TLLR/TESS $\to$ 阶段 5 逐层评估 $\to$ 阶段 6 硬件对接」，阶段 0/1 打勾、2--6 灰显。
\item \textbf{口播}：六阶段：0 搭骨架 $\to$ 1 跑穿 BPTT-SG $\to$ 2 焊钩子 $\to$ 3 挂 OTTT/NDOT/e-prop $\to$ 4 接 S-TLLR/TESS $\to$ 5 逐层评估 $\to$ 6 硬件对接；阶段 0/1 已完成，本次重点讲 0/1。
\end{itemize}

\textbf{第 5 页 · 阶段 0：把实验框架搭起来}

\begin{itemize}[leftmargin=1.5em,itemsep=0pt]
\item \textbf{幻灯片}：模块图 snn/（LIF 前向 + surrogate 反向）、network/（多层 SNN）、data/（泊松编码 + 合成数据）、utils/（训练循环 + 指标接口）、sanity/（验证）；右下角「全链路 PASS」。
\item \textbf{口播}：LIF（前向 + surrogate 反向）、多层 SNN、泊松编码 + 合成数据、训练循环、指标接口齐备；\texttt{sanity/stage0\_check.py} 全链路 PASS；意义：后续算法在同一框架跑、口径一致。
\end{itemize}

\textbf{第 6 页 · 阶段 1：手写 BPTT-SG，把反向跑穿}

\begin{itemize}[leftmargin=1.5em,itemsep=0pt]
\item \textbf{幻灯片}：反向递推公式——输出层 $\delta^L[t]=dL/ds_{\mathrm{out}}[t]\odot\Psi'+\lambda\,\delta^L[t{+}1]$；中间层 $\delta^l[t]=\big((W^{l+1})^\top\delta^{l+1}[t]\big)\odot\Psi'+\lambda\,\delta^l[t{+}1]$；权重梯度 $dW^l=\sum_t\delta^l[t]\otimes s^{l-1}[t]$；标注两种模式 detached-reset（$A_t=\lambda I$）与 full（$A_t=\lambda I-\lambda V_{th}D_\Psi$）。
\item \textbf{口播}：纯手写前向 + 反向递推，不依赖 autograd；同时实现 detached-reset（$A_t=\lambda I$）与 full（$A_t=\lambda I-\lambda V_{th}D_\Psi$）两模式；目的：阶段 3 挂 OTTT/NDOT 时能逐项对照。
\end{itemize}

\textbf{第 7 页 · 阶段 1 验证：Gate A/B/C 梯度对上了}

\begin{itemize}[leftmargin=1.5em,itemsep=0pt]
\item \textbf{幻灯片}：三行表——Gate A 时间 Jacobian $A_t$（detached $=\lambda I$ / full $=\lambda I-\lambda V_{th}D_\Psi$）；Gate B 空间 Jacobian $W D_\Psi$；Gate C 最终梯度 $dW/db$（多层多步）；右侧判据「手写 vs autograd 相对误差 $<10^{-5}$」，实测 $10^{-7}\sim10^{-8}$。
\item \textbf{口播}：三道 Gate：时间 Jacobian、空间 Jacobian、最终梯度 $dW/db$；手写 vs autograd 相对误差 $<10^{-5}$（实测 $10^{-7}\sim10^{-8}$）；结论：手写反向与自动微分逐项一致。
\end{itemize}

\textbf{第 8 页 · 阶段 1 收敛：手写梯度真的能训练}

\begin{itemize}[leftmargin=1.5em,itemsep=0pt]
\item \textbf{幻灯片}：full BPTT 收敛数据——网络 $8\text{-}64\text{-}3$、$T=10$、256 样本合成数据、Adam lr$=10^{-2}$、80 epochs，最后 3 个 epoch 准确率均值 $0.862$（方差 $1.5\times10^{-4}$，$>0.85$）；detached-reset 对照 $0.872$。
\item \textbf{口播}：设置：$8\text{-}64\text{-}3$、$T=10$、256 样本、Adam lr$=10^{-2}$、80 epochs；full BPTT last3 均值 $0.862$（$>0.85$），方差 $1.5\times10^{-4}$；结论：手写梯度能驱动训练收敛。
\end{itemize}

\textbf{第 9 页 · 一个伏笔：丢复位项的代价}

\begin{itemize}[leftmargin=1.5em,itemsep=0pt]
\item \textbf{幻灯片}：256 样本：full $0.862$ / detached $0.872$（接近）；64 样本（单次实验观察，未统计）：detached 降至 $\approx0.81$、低于 full——数据难分时丢复位项的代价才显现。
\item \textbf{口播}：detached-reset 丢复位项 $\to$ $A_t=\lambda I$，即 OTTT 的时间近似；256 样本两者接近（full $0.862$ / detached $0.872$）；单次实验观察（64 样本，未统计）：detached 降至约 $0.81$、低于 full——丢复位项的系统性代价；留作阶段 3 对比 OTTT 的参照。
\end{itemize}

\textbf{第 10 页 · 冻结的约定：让后面不再重推符号}

\begin{itemize}[leftmargin=1.5em,itemsep=0pt]
\item \textbf{幻灯片}：约定汇总简表——LIF 复位形式 A（$u_{t+1}=\lambda(u_t-V_{th}s_t)+I$）；超参 $\lambda=0.5$、$V_{th}=1.0$、$\beta=4.0$；符号对照 \texttt{SYMBOL\_MAPPING.md}（论文 $W^{l\leftarrow l-1}$ $\leftrightarrow$ 代码 \texttt{weights[l-1]}，差 1 层）。
\item \textbf{口播}：统一 LIF 复位形式 A（与 OTTT/NDOT/e-prop 符号最接近）；超参冻结：$\lambda=0.5$、$V_{th}=1.0$、$\beta=4.0$；符号对照 \texttt{SYMBOL\_MAPPING.md}：论文 $W^{l\leftarrow l-1}$ 对应代码 \texttt{weights[l-1]}（差 1 层）。
\end{itemize}

\textbf{第 11 页 · 下一步：阶段 2 起}

\begin{itemize}[leftmargin=1.5em,itemsep=0pt]
\item \textbf{幻灯片}：阶段 2--6 清单——阶段 2 焊探针/CKA/梯度对齐钩子（并在 BPTT 自身上验证 $\mathrm{CKA}=1$、$\cos=1$）；阶段 3 挂 OTTT/NDOT/e-prop；阶段 4 接 S-TLLR/TESS；阶段 5 逐层评估三张图 + 归因 + Pareto；阶段 6 硬件对接。
\item \textbf{口播}：阶段 2：焊探针/CKA/梯度对齐钩子，先在 BPTT 自身上验证 $\mathrm{CKA}=1$、$\cos=1$；阶段 3--5：挂各算法、逐层评估；阶段 6：整理成硬件可参考规格。
\end{itemize}

\textbf{第 12 页 · 需要讨论的问题}

\begin{itemize}[leftmargin=1.5em,itemsep=0pt]
\item \textbf{幻灯片}：4 个问题——(1) 逐层探针的准确率曲线是否真能定位信息损失（遗留疑问 H）；(2) 逐层 CKA 与梯度对齐是否一致（遗留疑问 I）；(3) 信息损失与资源节省的 Pareto 前沿形态（遗留疑问 J）；(4) 精读原文遗留疑问：OTTT Theorem 1 收敛条件、NDOT clamp 逻辑、DFA 收敛性证明、TESS $B$ 矩阵准正交性（遗留疑问 A/B/C/D）。
\item \textbf{口播}：探针曲线能否定位信息损失（H）；CKA 与梯度对齐结论是否一致（I）；Pareto 前沿形态（J）；精读原文遗留：OTTT Theorem 1、NDOT clamp、DFA 收敛证明、TESS $B$ 准正交。
\end{itemize}

\textbf{第 13 页 · 后续方向：基于逐层定位的算法改进}

\begin{itemize}[leftmargin=1.5em,itemsep=0pt]
\item \textbf{幻灯片}：逻辑链「逐层探针 / CKA / 梯度对齐 $\to$ 定位某层主要损失来源 $\to$ 针对该层改 OTTT/NDOT 的局部近似」，标注「后续方向 / 下一篇文章的切入点」。
\item \textbf{口播}：定位：非当前项目交付，后续文章种子；前提：阶段 2--5 逐层评估完成后才有依据；落点：探针/CKA/梯度对齐锁定某层损失来源，针对该层改 OTTT/NDOT 局部近似。
\end{itemize}

\subsection{2026-11-XX 组会}

% TODO：待补充。

\bibliographystyle{IEEEtran}
\bibliography{SNN_OnlineLearningSummary}
\end{document}


